\documentclass[11pt,a4paper]{article}

% ---- Fuentes Type1 reales (Times) + T1 encoding: evita tanto el bug de
% fuentes bitmap (ec/pk) como el bug de XeTeX que rompe los espacios entre
% palabras en el texto extraído de PDFs justificados. ----
\usepackage[T1]{fontenc}
\usepackage[utf8]{inputenc}
\usepackage{mathptmx}
% Nota: se omite babel[spanish] porque los patrones de guionizado en español
% no están instalados en este entorno; no afecta el parseo ATS del texto.

\usepackage[margin=1cm]{geometry}
\usepackage{titlesec}
\usepackage{enumitem}
\usepackage{hyperref}
\usepackage{xcolor}

\definecolor{darkgray}{RGB}{40,40,40}

\hypersetup{
    colorlinks=true,
    linkcolor=darkgray,
    urlcolor=darkgray
}

\usepackage{ragged2e}

\pagestyle{empty}
\setlength{\parindent}{0pt}
\setlength{\parskip}{2pt}
% Texto alineado a la izquierda (sin justificar): la justificación comprime
% el espacio entre palabras y eso es lo que rompe el parseo en ATS/pdftotext.
\RaggedRight

% Secciones simples, sin tablas ni multicolumna (clave para parseo ATS correcto)
% Nota: se evita \scshape (versalitas falsas) porque genera espacios fantasma
% en el texto extraído del PDF (ej. "P ERFIL P ROFESIONAL"), rompiendo la
% detección de encabezados de sección por parte de los ATS.
\titleformat{\section}{\large\bfseries}{}{0em}{}[\titlerule]
\titlespacing{\section}{0pt}{10pt}{6pt}

\newcommand{\job}[4]{%
    \textbf{\underline{#1}} \hfill #2\\
    \textit{#3} \hfill \textit{#4}\\[2pt]
}

\newcommand{\proj}[2]{%
    \textbf{#1} \\
    \textit{#2}\\[2pt]
}

\newcommand{\newsubsection}[1]{
    \vspace{0em}
    \textbf{\textit{#1}}
    \vspace{0em}
}


\begin{document}

% ---------- ENCABEZADO ----------
\begin{center}
    {\Huge \textbf{Jorge Pizarro Callejas}}\\[4pt]
    {\large Ingeniero de Software · Desarrollo Web \& Cloud}\\[6pt]
    +56 9 7624 7629 \; $\cdot$ \; \href{mailto:jpizarrocallejas@gmail.com}{jpizarrocallejas@gmail.com} \; $\cdot$ \;
    \href{https://linkedin.com/in/jorgepizarrocallejas}{linkedin.com/in/jorgepizarrocallejas} \; $\cdot$ \;
    \href{https://github.com/jorgicio}{github.com/jorgicio}\\[2pt]
    Quilpué, Chile
\end{center}

%\vspace{4pt}

% ---------- PERFIL PROFESIONAL ----------
\section*{Perfil Profesional}
Ingeniero de Ejecución en Informática titulado en la Universidad Técnica Federico Santa María, con más de 10 años de experiencia en desarrollo de software, especializado en aplicaciones web, servicios cloud (AWS) y automatización backend. Más de 3 años de experiencia comprobada en arquitecturas serverless con AWS Lambda, CI/CD con Jenkins y GitLab, y desarrollo fullstack con Java, Node.js, Angular y Vue.js. Orientado a la calidad del código, trabajo en equipo remoto y resolución de problemas técnicos complejos.\\
Actualmente busco oportunidades en el área de DevOps o desarrollo frontend para aportar mis conocimientos y aprender nuevas tecnologías en un entorno profesional desafiante.


% ---------- EXPERIENCIA LABORAL ----------
\section*{Experiencia Laboral}

\job{Ingeniero de Software Junior}{Octubre 2021 --- Enero 2025}{Tranciti (Waypoint Telecomunicaciones)}{Santiago, Chile}

\newsubsection{Funciones}
\begin{itemize}[leftmargin=*, noitemsep, topsep=2pt]
    \item Trabajé en el área de soporte de IT, gestionando la atención al cliente. 
    \item Desarrollé microservicios y funciones serverless con AWS Lambda en Java y Node.js para una plataforma de geolocalización GPS, mejorando la eficiencia en un 30\%.
    \item Implementé infraestructura como código (IaC) con AWS CloudFormation.
    \item Desarrollé un módulo frontend para clientes en Angular.
    \item Brindé soporte a apps móviles con Titanium, Ionic y Expo Go para 5 clientes.
    \item Utilicé PostgreSQL y AWS DynamoDB/Athena como motores de base de datos según el tipo de carga.
\end{itemize}

\newsubsection{Logros}
\begin{itemize}[leftmargin=*, noitemsep, topsep=2pt]
    \item Mejoré la satisfacción del cliente a través de la gestión del soporte.
    \item Reduje el tiempo de despliegue en un 40\% a través de la automatización con Jenkins.
    \item Reduje los errores de producción en un 25\% al integrar Docker y utilizar herramientas de calidad de código.
    \item Optimicé consultas y reduje hasta un 20\% el tiempo de respuesta en los motores de base de datos.
\end{itemize}

\job{Desarrollador Frontend}{Diciembre 2019}{Innovación y Proyectos Kimun Ltda.}{Temuco}

\newsubsection{Funciones}
\begin{itemize}[leftmargin=*, noitemsep, topsep=2pt]
    \item Desarrollé y realicé mejoras a nivel de frontend en la plataforma web ``Mundo de Vegetales'' utilizando HTML5, CSS3 y JavaScript.
\end{itemize}

\newsubsection{Logros}
\begin{itemize}[leftmargin=*, noitemsep, topsep=2pt]
    \item Reduje la carga del sitio hasta un 50\% y aumenté la retención de usuarios.
\end{itemize}

\job{Desarrollador WordPress}{Octubre 2016}{Serakin}{Viña del Mar}

\newsubsection{Funciones}
\begin{itemize}[leftmargin=*, noitemsep, topsep=2pt]
    \item Realicé el diseño y desarrollo de sitios web corporativos mediante la utilización de CMS como WordPress para la gestión y edición de los contenidos, la integración de elementos visuales y funcionalidades necesarias para el sitio web.
\end{itemize}

\newsubsection{Logros}
\begin{itemize}[leftmargin=*, noitemsep, topsep=2pt]
    \item Mejoré la visibilidad y la accesibilidad en línea del negocio.
    \item Facilité la actualización de la información del sitio.
\end{itemize}

\job{Desarrollador Frontend}{Octubre 2015 --- Diciembre 2015}{Everis}{Santiago}

\newsubsection{Funciones}
\begin{itemize}[leftmargin=*, noitemsep, topsep=2pt]
    \item Desarrollé el panel de clientes del sitio web de Consorcio usando Bootstrap y jQuery.
\end{itemize}
\newsubsection{Logros}
\begin{itemize}[leftmargin=*, noitemsep, topsep=2pt]
    \item Facilité la experiencia del usuario a través de un diseño más intuitivo y accesible.
\end{itemize}


\job{Desarrollador Web}{Septiembre 2014 --- Abril 2015}{Pontificia Universidad Católica de Valparaíso}{Viña del Mar}
\newsubsection{Funciones}
\begin{itemize}[leftmargin=*, noitemsep, topsep=2pt]
    \item Continué el desarrollo de un sistema de aprendizaje de inglés para un proyecto CONICYT usando PHP, jQuery y MySQL.
\end{itemize}

\newsubsection{Logros}
\begin{itemize}[leftmargin=*, noitemsep, topsep=2pt]
    \item Mejoré la experiencia de usuario en un 20\%.
\end{itemize}

\job{Desarrollador Web y Jefe de Área de Informática}{Junio 2011 --- Marzo 2014}{Del Puerto Comunicaciones}{Valparaíso}
\newsubsection{Funciones}
\begin{itemize}[leftmargin=*, noitemsep, topsep=2pt]
    \item Desarrollé y mantuve más de 5 sitios web para clientes de agencia usando WordPress, HTML, CSS y jQuery. Práctica profesional de 360 horas continuada como empleado de planta.
    
\end{itemize}
\newsubsection{Logros}
\begin{itemize}[leftmargin=*, noitemsep, topsep=2pt]
    \item Con mi incorporación a la empresa se creó el área de informática, la cual terminé liderando.
\end{itemize}

\job{Desarrollador Web}{Marzo 2013 --- Enero 2014}{Oxalix}{Viña del Mar}
\newsubsection{Funciones}
\begin{itemize}[leftmargin=*, noitemsep, topsep=2pt]
    \item Desarrollé formularios y módulos web con Python y Django, integrando MySQL.
\end{itemize}
\newsubsection{Logros}
\begin{itemize}[leftmargin=*, noitemsep, topsep=2pt]
    \item Logré crear una arquitectura de software más robusta y escalable para futuras actualizaciones.
\end{itemize}

% ---------- EDUCACIÓN ----------
\section*{Educación}
\textbf{Ingeniería de Ejecución en Informática} \hfill 2012 --- 2015\\
\textit{Universidad Técnica Federico Santa María} \hfill \textit{Valparaíso, Chile}\\
Titulado

% ---------- HABILIDADES TÉCNICAS ----------
\section*{Habilidades Técnicas}
\textbf{Lenguajes - Avanzado:} Java, JavaScript (Node.js), Python, PHP, C/C++, HTML5, CSS3, SQL. \textbf{Intermedio:} TypeScript.\\
\textbf{Frameworks y bibliotecas - Intermedio:} Angular, Vue.js, Django, Bootstrap, Tailwind CSS, Ionic, Expo Go, Quarkus.\\
\textbf{Cloud y DevOps - Intermedio:} AWS Lambda, AWS CloudFormation, AWS CodeCommit, AWS CloudWatch, AWS Glue, AWS Redshift, AWS QuickSight, Docker, Jenkins (CI/CD), SonarQube.\\
\textbf{Datos e Ingeniería de Datos - Avanzado:} PostgreSQL, MySQL.  \textbf{Intermedio:} AWS DynamoDB, AWS Athena, ETL, web scraping (BeautifulSoup), procesamiento de datos con Python.\\
\textbf{Herramientas - Avanzado:} Git, GitHub, GitLab, VS Code, Linux, \LaTeX.\\
\textbf{Metodologías:} CI/CD, control de versiones, QA automatizado, DevOps, Test-Driven Development (TDD), Domain-Driven Design (DDD).\\
%\textbf{Habilidades blandas:} Adaptabilidad, trabajo en equipo, resolución de problemas, orientación al cliente, comunicación efectiva

% ---------- PROYECTOS PERSONALES ----------
\section*{Proyectos Personales}

\begin{itemize}[leftmargin=*, noitemsep, topsep=2pt]
    \item \textbf{Scraper de ofertas laborales GetOnBoard} \textit{Hecho en Python con BeautifulSoup, PostgreSQL, Docker}.\\
    Pipeline de web scraping para extracción y almacenamiento estructurado de ofertas laborales, contenerizado con Docker y persistido en PostgreSQL.
    
    \item \textbf{Bot de Indicadores Económicos (Telegram)} \textit{Hecho en Node.js con API REST}.\\
    Bot en Telegram que responde comandos para mostrar indicadores económicos en tiempo real usando la API de mindicador.cl.
    \item \textbf{Sitio Web Personal con Portafolio} \textit{Hecho en Vue.js y Tailwind CSS, backend con GraphQL}.\\
    Sitio en desarrollo con carta de presentación, portafolio y blog.

    \item \textbf{Tomython -- Temporizador Pomodoro} \textit{Hecho en Python con PyGTK y GStreamer}.\\
    Notificador Pomodoro de escritorio con notificaciones nativas de Linux, usando PyGTK y GStreamer.
    
    \item \textbf{Blog Técnico en Inglés} \textit{Hecho en Python con Pelican}.\\
    Blog estático creado con el generador Pelican en Python, publicado en inglés.

\end{itemize}

\section*{Idiomas}
Español (nativo), Inglés (escrito: nivel avanzado, hablado: nivel intermedio).

\end{document}